% Generated draft: requires mathematical and source review.
\documentclass[11pt]{article}
\usepackage[margin=25mm]{geometry}
\usepackage{fontspec}
\setmainfont{Latin Modern Roman}
\usepackage{amsmath,amssymb,mathtools,unicode-math}
\setmathfont{Latin Modern Math}
\usepackage{enumitem,needspace,xurl,hyperref}
\setlength{\parindent}{0pt}
\setlength{\parskip}{5pt}
\setlength{\emergencystretch}{3em}
\newcommand{\N}{\mathbb{N}}
\newcommand{\Z}{\mathbb{Z}}
\newcommand{\Q}{\mathbb{Q}}
\newcommand{\R}{\mathbb{R}}
\newcommand{\C}{\mathbb{C}}
\newcommand{\F}{\mathbb{F}}
\newcommand{\A}{\mathbb{A}}
\newcommand{\PP}{\mathbb{P}}
\newcommand{\E}{\mathbb{E}}
\newcommand{\eps}{\varepsilon}
\newcommand{\dd}{\,\mathrm{d}}
\begin{document}
\section*{UnsolvedMath ID: 20003357}
\textbf{Source problem code:} catalogue entry
\subsection{Definitions and mathematical statement}
The complete local text is the incomplete bridge ``Conjecture 54 for $k\geq3$ follows from''. It contains neither Conjecture 54 nor the alleged antecedent and is not a standalone arithmeticity theorem. The aid recovers the continuation as a reference to Conjecture 29, but reports a dimension mismatch that must not be hidden.

Use $C_k$ only as a name for the relevant dimension-$k$ diagonal-orbit conjecture and $M_k$ for the dimension-$k$ arithmeticity assertion in the missing Conjecture 54 setup. The aid attributes to Hee Oh the implication
\[
 C_k\ \Longrightarrow\ M_k\qquad(k\geq3),
\]
whereas literal AIM Conjecture 29 states only the dimension-three case $C_3$. It does not follow from this citation alone that $C_3\Rightarrow M_k$ for every higher $k$. The actual groups, lattices and arithmeticity conclusion of $M_k$ are not reproduced in the supplied record, so the labels do not stand in for invented definitions.

Independently the aid reports that Benoist--Oh (2010) proves the intended Conjecture 54 for all $k\geq2$, and gives a second verification using Benoist--Miquel (2020). Its claimed hypothesis bridge concerns a discrete group with prescribed full lattices in opposite unipotent groups and projection-injectivity conditions. It says projection injectivity is exactly the irreducibility required \emph{in that precise AIM setup}. This is not a universal equivalence for every unrelated group product or every lattice in a vector space; the relevant projection targets and irreducibility definition are absent and need the full source.

The reported enlargement step has the following conditional content. If a lattice $\Lambda$ in a unipotent group $U$ is contained in the discrete intersection $\Gamma\cap U$, the latter is a finite-index superlattice. A projection-injectivity condition extends to it by taking a power into $\Lambda$ and using torsion-freeness of $U$. The aid uses this to match theorems stated for the full unipotent intersections rather than only the initially supplied lattices. The exact hypotheses ensuring discreteness and the prescribed projections must still be retained.

For the Zariski-density step, the aid reports that the full opposite unipotent groups in its setup are generated algebraically by the Zariski closures of their lattices and together generate the ambient group. Thus their inclusion forces the ambient subgroup to be Zariski dense, without an additional independent density assumption in that formulation. Opposite unipotent lattices alone are not silently declared an arithmeticity theorem in an arbitrary ambient group.

These are a dimension correction and an attributed unconditional-resolution/hypothesis audit for the missing neighboring problem. The present record supplies no full arithmeticity statement, complete proof or new independent conjecture. Its solved-in-literature label refers to the aid's reconstructed context and requires checking the actual cited results.
\subsection{Short English statement}
The record is an unfinished implication bridge. The aid distinguishes a dimension-matched historical implication from a stronger inference the fragment might suggest, and reports an independent arithmeticity resolution. The underlying groups and projection hypotheses are missing, so the equivalence claims are kept specific to the original setup.
\subsection{Sources}
\begin{itemize}
\item Local UnsolvedMath registry, permanent ID 20003357, source code catalogue entry; \texttt{aim-other-notes.json}, notes[163], source tag conjecture.
\item AIM, \emph{Emerging applications of measure rigidity}, fragment relating Conjecture 54 and Conjecture 29: \url{https://aimath.org/WWN/measrigid/measrigid.pdf}. Full adjacent statements are absent.
\item Hee Oh's dimension-matched implication, Benoist--Oh (2010), and Benoist--Miquel (2020), named in the supplied summary without titles, URLs or complete theorem statements.
\item Machine-generated research attempt 1: partial result, solved status at the run, rigor labelled ``proved,'' low novelty confidence; registry solved/``SOLVED-IN-LITERATURE.'' Not peer reviewed; independently verify before citation. Registry updated 14 August 2026.
\end{itemize}
\end{document}
